% ======================================================================
\chapter{Preparando o ambiente}
\label{cap:ambiente}
% ======================================================================

Duas ferramentas acompanham toda a capacitação:

\begin{itemize}
  \item o \textbf{Git}, que controla as versões no seu computador;
  \item o \textbf{GitHub CLI} (\texttt{gh}), que opera o GitHub pelo terminal:
        cria repositórios, abre Pull Requests e faz o login.
\end{itemize}

Siga só a seção do seu sistema.

% ----------------------------------------------------------------------
\section{Windows}
% ----------------------------------------------------------------------

No Windows, a instalação é feita pelo \textbf{winget}, o gerenciador de pacotes
da Microsoft \cite{microsoft_winget}. Ele vem no Windows 11 e nas versões
atuais do Windows 10 (a partir da 1809), como parte do aplicativo
\emph{App Installer}.

\subsection{Conferir o winget}

Abra o \textbf{PowerShell} (menu Iniciar \(\rightarrow\) digite
\enquote{PowerShell}) e rode:

\begin{bashcode}
winget --version
\end{bashcode}

Se aparecer um número de versão (Figura~\ref{fig:winget_installed}), está
pronto. Se o comando não for reconhecido, abra a Microsoft Store, procure
\enquote{App Installer} e atualize ou instale.

\begin{figure}[H]
  \centering
  \includegraphics[width=0.9\textwidth]{./assets/images/07_winget_installed.png}
  \caption{Conferindo o winget}
  \label{fig:winget_installed}
\end{figure}

\begin{dica}[Primeiro uso]
Na primeira vez, o winget pede para aceitar os termos de uso das fontes de
pacotes. Digite \texttt{Y} e aperte \tecla{Enter}.
\end{dica}

\subsection{Instalar o Git}

O \texttt{winget search} mostra o identificador (\emph{Id}) de cada programa
(Figura~\ref{fig:winget_search}). O Id do Git é \texttt{Git.Git}:

\begin{bashcode}
winget search git
winget install --id Git.Git -e --source winget
\end{bashcode}

\begin{figure}[H]
  \centering
  \includegraphics[width=0.9\textwidth]{./assets/images/08_winget_search.png}
  \caption{Buscando o Git no winget}
  \label{fig:winget_search}
\end{figure}

A opção \texttt{-e} exige o Id exato, e \texttt{--source winget} evita a
Microsoft Store. É a forma recomendada pelo site oficial do Git
\cite{git_windows}. Aguarde o fim da instalação
(Figura~\ref{fig:install_git}).

\begin{figure}[H]
  \centering
  \includegraphics[width=0.9\textwidth]{./assets/images/09_install_git.png}
  \caption{Instalação do Git}
  \label{fig:install_git}
\end{figure}

O Git para Windows já traz o \textbf{Git Bash}, o terminal que usaremos daqui
em diante, o \textbf{Git LFS} (capítulo~\ref{cap:lfs}) e o \textbf{Git
Credential Manager}, que guarda o seu login com segurança.

\subsection{Instalar o GitHub CLI}

\begin{bashcode}
winget install --id GitHub.cli
\end{bashcode}

\begin{figure}[H]
  \centering
  \includegraphics[width=0.9\textwidth]{./assets/images/11_install_cli.png}
  \caption{Instalando o GitHub CLI}
  \label{fig:install_cli}
\end{figure}

\subsection{Conferir}

\textbf{Feche o terminal e abra o Git Bash} (menu Iniciar \(\rightarrow\)
\enquote{Git Bash}). O terminal precisa ser reaberto para enxergar os
programas recém-instalados.

\begin{bashcode}
git --version
gh --version
\end{bashcode}

\begin{figure}[H]
  \centering
  \includegraphics[width=0.9\textwidth]{./assets/images/10_git_version.png}
  \caption{Versão do Git}
  \label{fig:git_version}
\end{figure}

\begin{figure}[H]
  \centering
  \includegraphics[width=0.9\textwidth]{./assets/images/12_cli_version.png}
  \caption{Versão do GitHub CLI}
  \label{fig:cli_version}
\end{figure}

\begin{dica}[Atualizar depois]
Para atualizar um programa específico: \texttt{winget upgrade --id Git.Git}.
Evite o \texttt{winget upgrade --all} no meio da aula: ele atualiza todos os
programas do computador e pode demorar bastante.
\end{dica}

% ----------------------------------------------------------------------
\section{Linux (Debian e Ubuntu)}
% ----------------------------------------------------------------------

O Git está nos repositórios da distribuição:

\begin{bashcode}
sudo apt update
sudo apt install git git-lfs
\end{bashcode}

O GitHub CLI tem repositório próprio. Este é o comando oficial
\cite{gh_linux}. Copie inteiro: as barras invertidas no fim das linhas unem
tudo num comando só.

\begin{bashcode}
(type -p wget >/dev/null || (sudo apt update && sudo apt install wget -y)) \
  && sudo mkdir -p -m 755 /etc/apt/keyrings \
  && out=$(mktemp) && wget -nv -O$out https://cli.github.com/packages/githubcli-archive-keyring.gpg \
  && cat $out | sudo tee /etc/apt/keyrings/githubcli-archive-keyring.gpg > /dev/null \
  && sudo chmod go+r /etc/apt/keyrings/githubcli-archive-keyring.gpg \
  && sudo mkdir -p -m 755 /etc/apt/sources.list.d \
  && echo "deb [arch=$(dpkg --print-architecture) signed-by=/etc/apt/keyrings/githubcli-archive-keyring.gpg] https://cli.github.com/packages stable main" | sudo tee /etc/apt/sources.list.d/github-cli.list > /dev/null \
  && sudo apt update \
  && sudo apt install gh -y
\end{bashcode}

\begin{aviso}[Evite o pacote gh da distribuição]
Os mantenedores do GitHub CLI desaconselham o \texttt{gh} empacotado pela
comunidade Debian: algumas versões antigas dele deixaram de funcionar com as
APIs atuais do GitHub \cite{gh_linux}. Use o repositório oficial acima.
\end{aviso}

Outras distribuições (Fedora, Arch, openSUSE) estão em
\url{https://github.com/cli/cli/blob/trunk/docs/install_linux.md}.

% ----------------------------------------------------------------------
\section{macOS}
% ----------------------------------------------------------------------

O caminho mais simples é o \textbf{Homebrew} (\url{https://brew.sh}). Com ele
instalado:

\begin{bashcode}
brew install git git-lfs gh
\end{bashcode}

\begin{dica}
Sem o Homebrew, rodar \texttt{git --version} no Terminal oferece instalar as
\emph{Command Line Tools} da Apple, que trazem o Git. O GitHub CLI, nesse
caso, é instalado pelo pacote em \url{https://cli.github.com}.
\end{dica}

% ----------------------------------------------------------------------
\section{Conferindo nos três sistemas}
% ----------------------------------------------------------------------

\begin{bashcode}
git --version
gh --version
git lfs version
\end{bashcode}

Os três devem responder com um número de versão. Se algum comando não for
reconhecido, feche e reabra o terminal. Se continuar, reinstale a ferramenta.
